\chapter{Conclusões e Trabalhos Futuros}
\label{cap:conclusao}

Neste trabalho, foi apresentada a \texttt{mlquantify}, uma biblioteca de quantificação em
Python construída sobre a interface do \textit{scikit-learn}. A motivação veio da
dispersão das implementações disponíveis, que obriga o usuário a adaptar dados e métodos
ao formato de cada pacote.

Na biblioteca, 45 métodos são reunidos sob uma interface única. Na arquitetura, são
combinados \textit{mixins} e \textit{tags}, o que permite acomodar exigências distintas
sem multiplicar classes base, e a decomposição multiclasse é aplicada de forma
automática a 19 desses métodos. Nove deles são composições do \textit{framework}
unificado, em que representação, função de perda e solucionador são escolhidos de
maneira independente. Sete métodos não estão disponíveis em nenhuma das bibliotecas de
referência, e apenas o PACC delas ficou sem implementação direta.

Além dos métodos, são disponibilizados quatro protocolos de avaliação, busca de
hiperparâmetros guiada pelo erro de quantificação e uma função em que divisão,
treinamento, amostragem e cálculo de métricas são executados em uma chamada. A
documentação é gerada a partir do próprio código, com 24 exemplos executáveis, e o
conjunto de testes é composto por 220 funções executadas a cada alteração do
repositório em cinco versões do Python.

Na avaliação experimental, a biblioteca proposta foi comparada com a \verb'QuaPy', a
\verb'QuantificationLib' e a \verb'qunfold' em 28 métodos e 132 conjuntos de dados, no
total de 39 pares de método e biblioteca, sempre sobre as mesmas amostras.

Quanto ao erro, no teste de equivalência o intervalo de confiança ficou dentro da margem
de um ponto percentual de prevalência em 38 dos 39 pares. Em 28 deles não há
diferença detectável, e nos 10 restantes a diferença é sistemática e irrelevante. A
exceção é o MS2 contra a \verb'QuaPy', favorável à biblioteca proposta e explicada pelo
conjunto de limiares candidatos de cada implementação. Fica assim demonstrada a
equivalência, afirmação mais forte que a simples ausência de diferença.

Quanto ao custo, a razão mediana ficou em $0{,}967$ no tempo e em $0{,}816$ na memória. A
vantagem em memória é ampla, com 32 dos 39 pares abaixo da igualdade. Em tempo há
equilíbrio, com perdas concentradas na contagem ajustada da \verb'QuantificationLib',
que resolve o ajuste em forma fechada, e ganhos nos métodos multiclasse de casamento de
distribuição, onde a \verb'qunfold' recorre a otimização numérica a cada estimativa.

O objetivo do trabalho fica sustentado por esses resultados. Entre as quatro bibliotecas
comparadas, a proposta é a de cobertura mais ampla de métodos, com erro equivalente ao
das implementações de referência e custo da mesma ordem de grandeza, sob uma interface
que dispensa adaptação de dados e de métodos.

O alcance dessas conclusões é delimitado por algumas limitações. O classificador base foi
o mesmo em toda a comparação, escolha que isola o quantificador e deixa em aberto o
efeito do modelo subjacente. Os hiperparâmetros foram os padrões de cada biblioteca, de
modo que o resultado descreve o uso sem ajuste. A margem de equivalência está na unidade
da prevalência e vale para o AE, o que exige nova escolha de margem para outras medidas
de erro. As medidas de custo vieram de uma única máquina, em processo único, e descrevem
razões entre bibliotecas em vez de tempos absolutos. Por fim, os métodos baseados em
máquinas de vetores de suporte e em redes neurais ficaram fora da comparação.

Como continuidade, cinco frentes se destacam. A primeira é expor a varredura exata dos
escores na família de limiares, hoje fixada em uma grade de 101 pontos, única origem de
divergência fora da margem de equivalência. A segunda é adotar solução fechada nos
métodos de contagem ajustada, onde está concentrada a desvantagem de tempo.

A terceira frente é a medida do esforço de uso, prevista no método e deixada para
avaliação posterior, com a contagem de linhas de código e de configurações manuais
necessárias para conduzir um mesmo experimento em cada biblioteca. A quarta é ampliar a
comparação para outros classificadores base e para cenários com ajuste de
hiperparâmetros, o que aproximaria a avaliação do uso em pesquisa. A quinta é incorporar
à comparação os métodos baseados em máquinas de vetores de suporte e em redes neurais,
já implementados na biblioteca.

A \texttt{mlquantify} está publicada no PyPI e, no repositório, estão o código, a
documentação e os \textit{scripts} dos experimentos, de modo que os resultados deste
trabalho podem ser reproduzidos e estendidos.
